\begin{figure}[htbp]
  \centering
  \resizebox{0.92\textwidth}{!}{%
  \begin{tikzpicture}[
    >=Latex,
    font=\small,
    node distance=12mm and 14mm,
    classbox/.style={
      draw,
      rounded corners=3pt,
      minimum width=58mm,
      align=left,
      fill=blue!6
    },
    objectbox/.style={
      draw,
      rounded corners=3pt,
      minimum width=58mm,
      align=left,
      fill=green!7
    },
    arrow/.style={-{Latex[length=2.3mm]}, thick}
  ]
    \node[classbox] (class) {
      \textbf{Class: Melody}\\[1mm]
      Attributes:\\
      \quad title\\
      \quad events\\
      \quad key\\[1mm]
      Methods:\\
      \quad add\_event()\\
      \quad get\_length()
    };

    \node[objectbox, right=28mm of class] (object) {
      \textbf{Object: melody\_a}\\[1mm]
      title = ``Example theme''\\
      events = (c4, d4, e4, g4, \ldots)\\
      key = C major
    };

    \draw[arrow] (class) -- node[above, align=center] {instantiated as\\a concrete melody} (object);
  \end{tikzpicture}%
  }
  \caption{A simple illustration of the difference between a class and an object. A class describes what kind of data and behaviour an object has, while an object stores one concrete instance.}
  \label{fig:oop_class_object}
\end{figure}
